\section{Study 3: a guaranteed false-alarm rate}\label{sec:conf}
The Bayesian $h_0$ makes few confident errors, but it promises nothing about how often it cries ``unknown'' by mistake. So we added a second decision rule that does, and we tested it on supernovae nobody had used before.

\paragraph{Method} After any planner's measurements $D$, we compute $s=\log p(D\mid h_0)-\max_c\log p(D\mid c)$ and flag ``unknown'' when $s$ goes above a threshold $\tau$. We set $\tau$ by split-conformal calibration on known supernovae, so if new supernovae look like the calibration ones, the chance of a false ``unknown'' is at most $\alpha=\ConfAlpha$, for any planner \citep{Vovk2005,Angelopoulos2023}. We also added a lookahead planner L, which gives some value ($\lambda=\ConfLambda$, chosen on validation data) to keeping later epochs available, because the greedy planner used too little of its budget. We drew \ConfNCal\ calibration and \ConfNFresh\ sealed test supernovae from Bright Transient Survey events that no earlier analysis had touched.

\begin{table*}[t]
\centering\scriptsize
\caption{Conformal open-set study, preregistered results per noise seed: estimate [95\% CI]; $^\times$ marks a criterion not met. Primary system: lookahead planner L with the conformal decision (C-L). $\Delta$: C-L minus closed-world EIG (B3), except HC6 (AUROC of $s$, L minus P). HC1 uses Clopper--Pearson intervals; HC2--HC4 and HC6 use the paired bootstrap.}
\label{tab:conf}
\begin{tabular}{llll}
\toprule
Hypothesis & Seed 1 & Seed 2 & Seed 3 \\
\midrule
HC1 false ``unknown'', sealed SNe ($\leq\alpha$) & +0.035 [+0.020, +0.055] & +0.061 [+0.042, +0.086] & +0.063 [+0.043, +0.089] \\
HC2 $\Delta$acc.\ SNe ($>-0.05$) & +0.082 [+0.039, +0.123] & +0.074 [+0.033, +0.115] & +0.059 [+0.012, +0.106] \\
HC3 $\Delta$CW SNe ($<+0.05$) & -0.027 [-0.057, +0.002] & -0.049 [-0.078, -0.020] & -0.012 [-0.043, +0.018] \\
HC4 $\Delta$CW unknowns ($<0$) & -0.239 [-0.375, -0.102] & -0.261 [-0.409, -0.114] & -0.193 [-0.341, -0.045] \\
HC5 unknowns flagged ($>\alpha$) & +0.420 [+0.316, +0.530] & +0.545 [+0.436, +0.652] & +0.432 [+0.327, +0.542] \\
HC6 $\Delta$AUROC, L$-$P ($>0$) & +0.012 [-0.013, +0.038]$^\times$ & +0.031 [+0.006, +0.058] & +0.039 [+0.005, +0.075] \\
\bottomrule
\end{tabular}

\end{table*}

\paragraph{The guarantee held, and accuracy went up} On the sealed supernovae, the false-``unknown'' rate was \ConfFa, inside the promised \ConfAlphaPct. The lookahead planner with the conformal rule was also \ConfAccGain\ more accurate than closed-world information gain, which settles the cost question that Study~2 left open. On the real unknowns it flagged \ConfFlag. Lookahead did not, however, rank the unknowns any better than P did (Table~\ref{tab:conf}).

\paragraph{Two rules, two kinds of safety} The conformal rule works on top of any planner, even a closed-world one, which then flagged \ConfFlagBthreeConf\ of the unknowns. But it is more often confidently wrong on unknowns (\ConfCwUnkCL\ against \ConfCwUnkP\ for P's Bayesian rule). So the conformal rule buys a certified false-alarm rate and accuracy on known classes, and the Bayesian $h_0$ buys fewer confident errors on the unknown.

\paragraph{A guarantee with no calibration at all} The conformal rule still needs calibration events, so we also tried the e-process of Eq.~\eqref{eq:eprocess}, which needs none. We replayed every planner exactly, measurement for measurement, and raised ``unknown'' the first time $E_t$ reached 10. Real light curves never match a template exactly, so this was a real test of the guarantee. On the \EFreshN\ sealed supernovae, P's false-``unknown'' rate was only \EFreshFa\ against the promised 10\%, and it still caught \EDet\ of the \EWithheldN\ real unknowns. That is fewer than $P(h_0)$ catches (\RealCFlagP), but the guarantee now holds at every step. Because the test works on any planner, it caught even more on planners that measure more (\EDetL\ for lookahead), so a planner that ignores $h_0$ can still be given a safe ``unknown'' label afterwards.
